\item[]\textbf{Unidad IV. Precipitación}\par\medskip

\question [5] Una lluvia de origen orográfico se favorece cuando:
\begin{choices}
\choice una masa de aire desciende por sotavento y se calienta.
\CorrectChoice una masa de aire húmedo asciende forzada por una barrera topográfica y se enfría.
\choice el calentamiento superficial produce convección libre sin relieve.
\choice un frente frío no altera el movimiento vertical del aire.
\end{choices}

\question [5] La diferencia más importante entre pluviómetro y pluviógrafo es que:
\begin{choices}
\choice ambos miden exactamente lo mismo y entregan el mismo tipo de registro.
\CorrectChoice el pluviómetro registra lluvia acumulada y el pluviógrafo permite observar su variación en el tiempo.
\choice el pluviógrafo solo se usa para nieve y el pluviómetro solo para lluvia líquida.
\choice el pluviómetro siempre entrega intensidades instantáneas y el pluviógrafo acumulados diarios.
\end{choices}

\question [5] En cuatro estaciones se midieron \SI{42}{mm}, \SI{38}{mm}, \SI{46}{mm} y \SI{34}{mm}. Si se usa el promedio aritmético para estimar la precipitación media areal, el valor es:
\begin{choices}
\choice \SI{38.0}{mm}
\choice \SI{39.0}{mm}
\CorrectChoice \SI{40.0}{mm}
\choice \SI{41.5}{mm}
\end{choices}

\question [5] El método de polígonos de Thiessen es más adecuado que el promedio aritmético cuando:
\begin{choices}
\choice las intensidades son idénticas en toda la cuenca por definición.
\CorrectChoice las estaciones no están distribuidas uniformemente y se desea ponderar por área de influencia.
\choice solo existe una estación y no se conoce su ubicación exacta.
\choice la lluvia debe expresarse en unidades de caudal.
\end{choices}

\question [5] Tres estaciones vecinas registran para un evento \SI{78}{mm}, \SI{84}{mm} y \SI{90}{mm}. Sus precipitaciones normales son \SI{100}{mm}, \SI{120}{mm} y \SI{150}{mm}, respectivamente. Si la estación faltante tiene precipitación normal de \SI{140}{mm}, la lluvia estimada por razón normal es aproximadamente:
\begin{choices}
\choice \SI{82.0}{mm}
\choice \SI{88.0}{mm}
\CorrectChoice \SI{102.7}{mm}
\choice \SI{118.0}{mm}
\end{choices}

\question [4] La curva de doble masa se utiliza principalmente para:
\begin{choices}
\choice transformar precipitación en escurrimiento directo.
\choice estimar la infiltración acumulada durante una tormenta.
\CorrectChoice evaluar consistencia y detectar cambios sistemáticos en una serie de datos.
\choice convertir una lluvia puntual en período de retorno.
\end{choices}

\question [4] Una curva altura--área--duración muestra, en general, que:
\begin{choices}
\choice la altura media siempre aumenta cuando el área considerada aumenta.
\choice la duración carece de efecto sobre la lluvia media.
\CorrectChoice para una misma duración, la precipitación media areal tiende a disminuir al aumentar el área.
\choice las curvas solo sirven para evaporación y no para precipitación.
\end{choices}

\item[]\textbf{Unidad V. Escurrimiento}\par\medskip

\question [5] ¿Cuál combinación describe mejor el escurrimiento directo durante una tormenta?
\begin{choices}
\choice Solo el flujo base sostenido por almacenamiento subterráneo profundo.
\CorrectChoice La parte del hidrograma asociada a la respuesta relativamente rápida del evento lluvioso.
\choice Únicamente el agua que circula en tuberías urbanas.
\choice Todo el caudal medido, sin distinguir componentes.
\end{choices}

\question [5] En una sección de aforo, el área hidráulica es \SI{6.0}{m^2} y la velocidad media es \SI{1.8}{m/s}. El caudal estimado es:
\begin{choices}
\choice \SI{3.33}{m^3/s}
\choice \SI{7.80}{m^3/s}
\CorrectChoice \SI{10.8}{m^3/s}
\choice \SI{13.0}{m^3/s}
\end{choices}

\question [5] En un hidrograma de tormenta, el gasto pico corresponde:
\begin{choices}
\choice al inicio del ascenso del caudal.
\choice al punto donde termina el escurrimiento base.
\CorrectChoice al valor máximo de descarga alcanzado durante el evento.
\choice al caudal medio anual de la corriente.
\end{choices}

\question [5] En igualdad de lluvia y área drenada, la urbanización tiende a producir:
\begin{choices}
\choice menor gasto pico y mayor tiempo al pico por incremento de infiltración.
\CorrectChoice mayor gasto pico y respuesta más rápida por aumento de superficies impermeables.
\choice el mismo hidrograma si el cauce principal no cambia.
\choice mayor recesión base sin afectar el resto del hidrograma.
\end{choices}

\question [5] Un hidrograma unitario se define como:
\begin{choices}
\choice el hidrograma de cualquier tormenta con duración unitaria sin importar el exceso de lluvia.
\CorrectChoice el hidrograma de escurrimiento directo producido por una unidad de lluvia efectiva uniformemente distribuida sobre la cuenca en una duración dada.
\choice la curva de gastos asociada a un cauce aforado una sola vez.
\choice la porción del hidrograma correspondiente al flujo base.
\end{choices}

\question [4] La curva de gastos relaciona de manera principal:
\begin{choices}
\choice área de cuenca y precipitación media.
\CorrectChoice tirante o nivel de agua con caudal descargado en una sección.
\choice velocidad del viento con evaporación potencial.
\choice período de retorno con intensidad de lluvia.
\end{choices}

\question [4] Si dos cuencas reciben tormentas comparables, la que posee drenaje más eficiente y menor almacenamiento superficial tenderá a mostrar:
\begin{choices}
\choice recesión más lenta y pico menor.
\choice igual hidrograma, porque solo importa el total de lluvia.
\CorrectChoice tiempo al pico menor y concentración más rápida del escurrimiento.
\choice desaparición del flujo base permanente.
\end{choices}

\item[]\textbf{Unidad VI. Predicción de eventos de diseño}\par\medskip

\question [5] Un período de retorno de 25 años para una lluvia o caudal extremo significa que:
\begin{choices}
\choice el evento ocurre exactamente una vez cada 25 años.
\CorrectChoice la probabilidad anual de excedencia es de 4\% bajo las hipótesis del análisis.
\choice después de que ocurre, no puede repetirse antes de 25 años.
\choice la serie observada tiene obligatoriamente 25 datos.
\end{choices}

\question [5] Si el período de retorno de un evento es 50 años, su probabilidad anual aproximada de excedencia es:
\begin{choices}
\choice 0.5\%
\choice 1\%
\CorrectChoice 2\%
\choice 5\%
\end{choices}

\question [5] Para un evento con período de retorno de 50 años, la probabilidad de que ocurra al menos una vez en 25 años es aproximadamente:
\begin{choices}
\choice 20\%
\choice 33\%
\CorrectChoice 39.7\%
\choice 63.2\%
\end{choices}

\question [5] Antes de ajustar una distribución probabilística a máximos anuales, conviene verificar al menos que la muestra sea:
\begin{choices}
\choice periódica y perfectamente simétrica.
\CorrectChoice razonablemente homogénea e independiente.
\choice exclusiva de años húmedos.
\choice ordenada de menor a mayor caudal solamente.
\end{choices}

\question [5] En hidrología de frecuencia, EV1 de Gumbel y Log-Pearson Tipo III se mencionan con frecuencia porque:
\begin{choices}
\choice son métodos para medir velocidad en una sección aforada.
\choice reemplazan la necesidad de registros observados.
\CorrectChoice son modelos probabilísticos usados para describir variables extremas hidrológicas.
\choice se usan únicamente para evaporación diaria.
\end{choices}

\question [5] Una serie de máximos anuales contiene 20 datos. Si el mayor evento ocupa rango \(m=2\) al ordenar de mayor a menor, el período de retorno empírico por \(T = (n+1)/m\) es:
\begin{choices}
\choice 5.0 años
\CorrectChoice 10.5 años
\choice 19.0 años
\choice 21.0 años
\end{choices}

\question [4] En una curva IDF, para un mismo período de retorno, al aumentar la duración de la tormenta suele ocurrir que:
\begin{choices}
\choice la intensidad aumenta sin límite.
\choice la intensidad permanece constante por definición.
\CorrectChoice la intensidad disminuye, aunque la lámina acumulada pueda aumentar.
\choice desaparece la relación con el diseño hidrológico.
\end{choices}
